\chapter{Conclusion}
\label{chap:conclusion}

This project built a reproducible benchmarking framework for Monte Carlo option-pricing workloads with automatic differentiation. The framework supports typed configurations, a common engine abstraction, repeated timing, deterministic seeding, correctness validation, SQLite storage, cloud metadata, cost estimation and report artefact generation. It evaluates NumPy, JAX, C++ and Rust implementations across European, Asian and local-volatility workloads where supported.

The strongest result is the SHA cloud profiler. In the final \code{sha_cloud_profiler_v4} experiment on \code{n2-standard-4}, the full grid required 48 full benchmark runs. SHA selected 24, saving 50.0\% of full runs while preserving 100.0\% Pareto recovery and 0.0\% runtime and cost regret relative to the full grid. The old profiler also preserved the best selections, but selected 30 full runs, so SHA achieved the same measured selection quality with fewer full-scale evaluations.

The engine results show that implementation choice is workload dependent. Rust/Rayon was the fastest measured no-AD engine for the European and local-volatility rows highlighted in the final cloud experiment, while JAX was essential for AD and was the best no-AD engine for the Asian cloud rows. JAX AD overhead at \(M=100{,}000\) ranged from 1.81\(\times\) to 4.42\(\times\), depending on workload and AD mode. These results support the original motivation: MC+AD performance cannot be inferred from the programming model alone.

The main limitations are the restricted hardware and software scope. GPU execution, Mojo, wider BLAS comparisons, hardware counters, more cloud regions and distributed execution were not completed. Future work should extend the framework in those directions, increase the number of cloud repetitions, and test SHA on larger and more heterogeneous configuration grids.

Overall, the project contributes both an engineering artefact and an empirical methodology. The artefact makes MC+AD experiments reproducible and queryable; the methodology shows that cheap probes and SHA can substantially reduce full cloud benchmarking work while preserving selection quality for the measured grid.
